\section{Methods}
\label{sec:methods}

\subsection{Datasets}
\label{sec:methods_data}
\textbf{UCSD Stress} (primary). Resting-state EEG from
\citep{komarov2020stress}: 17 graduate students, 92 recordings.
Each recording carries two behavioural labels supplied by the source
release: a categorical DASS-21 \citep{lovibond1995} state assignment
(\textit{increase} or \textit{normal}, per recording) used as the
primary binary target, and a continuous Daily Stress Scale (DSS)
self-report score (normalised to $[0, 100]$). We use the 70
recordings with both labels available; epochs are $5$\,s at
$200$\,Hz $\times$ $30$ channels. The primary binary label
(\textit{increase} vs \textit{normal}) is taken directly from the
source \texttt{Group} field and matches the label used by Wang et
al.\ \citep{wang2025stress} on this cohort. DASS-21 is a past-week
trait/state screener and DSS a daily-stress aggregate; neither was
psychometrically validated as a per-recording single-session EEG
classification label, and we reproduce both label conventions for
independent audit rather than endorsing either as a per-recording
neural ground truth.
\label{sec:methods_dass_scope}

\emph{Sample-size conventions.} Three distinct subsets are used
across the paper, and figure/table captions name the subset used
locally:
\begin{itemize}
  \item $70$ recordings / $17$ subjects --- subject-level
    fine-tuning and cross-validation against the DASS-21 binary
    target (Secs.~\ref{sec:results_paired},
    \ref{sec:results_fm_adds}; Fig.~\ref{fig:cv_gap}).
  \item $55$ recordings / $14$ subjects --- variance decomposition
    of frozen embeddings (Sec.~\ref{sec:results_atlas}). The nested
    decomposition used there requires each subject to carry a
    single target label, so the three DASS-mixed subjects (pids
    with recordings in both the \textit{increase} and
    \textit{normal} classes, $15$ recordings in total) are removed
    for this analysis only.
  \item $14$ subjects with $\geq\!2$ recordings --- within-subject
    longitudinal DSS trajectory classification
    (Sec.~\ref{sec:results_paired} Arm 2), where each recording is
    labelled above vs below that subject's personal-median DSS.
\end{itemize}

\textbf{EEGMAT} \citep{eegmat} (paired within-subject comparator).
36 healthy subjects $\times$ 2 recordings (resting baseline plus
serial-subtraction mental arithmetic); 19 channels, 60\,s segments.
This dataset provides the \emph{strong within-subject
alpha-desynchronisation contrast} side of our paired experiment
(Sec.~\ref{sec:results_paired}).

\textbf{ADFTD} \citep{adftd} (for variance-decomposition
cross-dataset coverage). Alzheimer's disease vs healthy controls,
65 subjects $\times$ 19 channels. ADFTD is natively
single-recording-per-subject; for the nested variance decomposition
we require multiple observations per subject, so each subject's
$\sim\!14$-min resting recording is partitioned into $3$
non-overlapping contiguous segments (195 pseudo-observations from 65
subjects, all three segments of a subject carrying the same AD/HC
label, kept in the same fold by subject-level CV). Full rationale
and the semantic-asymmetry caveat vs other datasets are in
App.~\ref{app:adftd_split}.

\textbf{TDBRAIN} \citep{tdbrain} (for variance-decomposition
cross-dataset coverage). MDD vs HC, 734 recordings / 359 subjects /
26 channels (downsampled to 19 for FM compatibility).

ADFTD is used in this paper for frozen-representation variance
decomposition (Sec.~\ref{sec:results_atlas}) and as a single
suggestive-but-not-controlled cross-subject anchored datapoint
(Sec.~\ref{sec:results_fm_adds}). TDBRAIN is used for variance
decomposition and additionally as the single
intermediate-anchor case study discussed in
Sec.~\ref{sec:discussion}, where its role in refining the binary
anchored/bounded SDL classification is explained.

\subsection{Pipeline and foundation models}
\label{sec:methods_pipeline}
All three FMs are wrapped in a unified frozen-extractor interface
producing fixed-dimensional embeddings from (channels $\times$ time)
windows: LaBraM \citep{labram2024} 200-d, CBraMod
\citep{cbramod2025} 200-d, REVE \citep{reve2025} 512-d. Weights are
loaded from the official pretrained releases. Input normalisation is
model-dependent to match each FM's pretraining distribution:
$z$-score for LaBraM, $\mu$V-scale unchanged for CBraMod and REVE.

\subsection{Cross-validation and evaluation protocols}
\label{sec:methods_cv}
Primary evaluation is \textbf{subject-level StratifiedGroupKFold(5)},
grouping by \texttt{patient\_id}. Trial-level StratifiedKFold(5) is
included only as a literature-comparison reference (e.g., reproducing
Wang et al.\ \citep{wang2025stress}) and reported transparently
alongside subject-level numbers; it carries subject leakage and
inflates reported BA on this dataset by $20$--$30$\,pp
(Sec.~\ref{sec:results_paired}).

All training uses cuDNN determinism. Seed variance on the UCSD Stress
$70$-recording / $14$-positive regime nevertheless remains
$\pm 5$--$10$\,pp single-seed, consistent with recent small-$N$
stability literature \citep{smalln2025biorxiv}. All headline claims
are reported as mean $\pm$ sample standard deviation (divide by
$n-1$) over at least three random seeds.

\textbf{Permutation null.} Recording-level label shuffle prior to CV.
A batched helper runs ten permutation seeds per FM; real FT and null
FT are compared by a one-sided Mann--Whitney test on the pooled
balanced-accuracy distributions.

\textbf{Class imbalance handling for classical baselines.} All
classical baselines use class-balanced weighting: Random Forest,
Logistic Regression ($L_2$), and SVM via
\texttt{class\_weight = `balanced'}; gradient-boosted trees via
\texttt{sample\_weight} from
\texttt{sklearn.utils.compute\_sample\_weight}.

\subsection{Variance decomposition and representational similarity}
\label{sec:methods_variance}
For the variance decomposition
(Sec.~\ref{sec:results_atlas}, Fig.~\ref{fig:variance_atlas}) we
compute three measurements on the frozen embedding matrix
$\mathbf{X} = (\mathbf{x}_1, \ldots, \mathbf{x}_N) \in
\mathbb{R}^{N \times D}$ for each cell, with per-recording subject
identifiers $s_i$ and target labels $y_i$.

\textbf{(i) Pooled variance fractions.} Let $\bar{\mathbf{x}}$
denote the grand mean. The total sum of squares is
$\mathrm{SS}_\text{total} = \sum_i \| \mathbf{x}_i -
\bar{\mathbf{x}} \|^2$. For datasets where each subject in the
analyzed subset carries a single target label --- ADFTD, TDBRAIN,
and the $55$-recording Stress subset after dropping DASS-mixed
subjects (Sec.~\ref{sec:methods_data}) --- we use a nested
decomposition (subject within label):
\begin{align*}
\mathrm{SS}_\text{label}      &= \sum_{\ell} n_\ell \, \| \bar{\mathbf{x}}_\ell - \bar{\mathbf{x}} \|^2, \\
\mathrm{SS}_\text{subj|label} &= \sum_{\ell} \sum_{s \in \ell} n_{s,\ell} \, \| \bar{\mathbf{x}}_{s,\ell} - \bar{\mathbf{x}}_\ell \|^2, \\
\mathrm{SS}_\text{residual}   &= \sum_i \| \mathbf{x}_i - \bar{\mathbf{x}}_{s_i, y_i} \|^2,
\end{align*}
where $\bar{\mathbf{x}}_\ell$ and $\bar{\mathbf{x}}_{s,\ell}$ are
the label-group and subject-within-label means, and $n_\ell$,
$n_{s,\ell}$ the corresponding counts. For EEGMAT, each subject
contributes recordings of both labels (rest and arithmetic), so the
design is crossed subject $\times$ label:
\begin{align*}
\mathrm{SS}_\text{label}    &= \sum_{\ell} n_\ell \, \| \bar{\mathbf{x}}_\ell - \bar{\mathbf{x}} \|^2, \\
\mathrm{SS}_\text{subject}  &= \sum_{s} n_s \, \| \bar{\mathbf{x}}_s - \bar{\mathbf{x}} \|^2, \\
\mathrm{SS}_\text{residual} &= \mathrm{SS}_\text{total} - \mathrm{SS}_\text{label} - \mathrm{SS}_\text{subject}.
\end{align*}
Pooled variance fractions are reported as
$\mathrm{SS}_\text{factor} / \mathrm{SS}_\text{total}$, with the
squared norms summed over all $D$ feature dimensions before
ratio-taking.

\textbf{(ii) RSA} \citep{kriegeskorte2008rsa}. The feature
dissimilarity matrix uses cosine distance,
$\mathrm{RDM}_\text{feat}[i,j] = 1 - \cos(\mathbf{x}_i,
\mathbf{x}_j)$, while the identity matrices are binary,
$\mathrm{RDM}_\text{label}[i,j] = \mathbb{1}[y_i \neq y_j]$ and
$\mathrm{RDM}_\text{subj}[i,j] = \mathbb{1}[s_i \neq s_j]$. We
report Spearman rank correlations between the feature RDM and the
subject and label RDMs on the upper-triangular off-diagonal
entries.

\textbf{(iii) Cluster bootstrap}
\citep{field2007clusterbootstrap}. A $95$\% confidence interval on
the subject/label variance ratio is obtained by resampling subjects
(not recordings) with replacement for $B = 1000$ iterations; each
iteration concatenates the drawn subjects' recordings and recomputes
$\mathrm{SS}_\text{subject} / \mathrm{SS}_\text{label}$ under
whichever decomposition applies. The ratio statistic is
heavy-tailed, so we report percentile bounds on the
log-transformed distribution and exponentiate them back for
reporting.

\subsection{Contrast-strength anchoring protocol}
\label{sec:methods_anchor}
A naïve operationalisation of the SDL thesis would define ``contrast
strength'' in terms of FM fine-tuning outcomes and then use those
outcomes to validate the thesis --- a circular construction. We
avoid this by triangulating contrast strength from three
FM-independent dimensions, assessed per dataset prior to any FM
fine-tuning.

\textbf{Dimension 1 --- Published within-subject correlate.}
Literature search for a published per-recording within-subject EEG
correlate of the target label, with documented spectral locus and
effect size.

\textbf{Dimension 2 --- Empirical cluster-corrected contrast.}
Group-level cluster-permutation $F$-test on log-PSDs (Welch, 2\,s
segments, 50\% overlap, 0.5--45\,Hz, MNE
\texttt{permutation\_cluster\_test}, 1000 permutations, default
threshold at $p = 0.05$, independent adjacency across channels and
frequency bins) on our own cohort, testing whether any cluster
survives correction at $p < 0.05$.

\textbf{Dimension 3 --- Locus concordance.}
If an empirical cluster survives, does it spatially and spectrally
align with the published within-subject correlate from Dimension 1?
An empirical cluster at an unexpected locus (different band,
different topography) is diagnostically weaker than one matching the
published anchor, because it may reflect a cohort-level confound
rather than a canonical label-relevant within-subject contrast.

A dataset is classified as \emph{Strong-anchored} when all three
dimensions agree (published correlate, significant cluster, and
locus match), \emph{Intermediate} when the empirical cluster
survives but either the published correlate is absent or its
spectral locus does not match, and \emph{Bounded} when no
within-subject correlate is documented and our own cluster test
does not survive correction. Fig.~\ref{fig:sdl_spectrum} uses this
rubric to order datasets. Full cluster-permutation protocols and
topographic summaries for all three datasets we anchor quantitatively
(UCSD Stress, EEGMAT, TDBRAIN) are in App.~A.

The rationale is to falsify SDL if an anchor disagrees with the
behavioural pattern, rather than merely confirm it post hoc.

\subsection{Hyperparameter selection}
\label{sec:methods_hp}
Hyperparameters for each FM's best fine-tuning configuration on the
primary UCSD Stress cohort were selected via a $3 \times 2 \times 3$
grid (3 learning rates $\times$ 2 encoder-learning-rate scales $\times$
3 seeds) independently per model. The selected configurations differ
between models: LaBraM $\text{lr} = 10^{-4}$, $\text{elrs} = 1.0$;
CBraMod $\text{lr} = 10^{-5}$, $\text{elrs} = 0.1$; REVE
$\text{lr} = 3 \times 10^{-5}$, $\text{elrs} = 0.1$. On the three
secondary datasets (EEGMAT, ADFTD, TDBRAIN) we did not re-sweep the
full grid; instead we used a per-model conservative fine-tuning
configuration identical across secondary datasets to control for
hyperparameter choice as a source of cross-dataset variation
(CBraMod and REVE at their Stress-sweep best HP; LaBraM at
$\text{lr} = 10^{-5}$, $\text{elrs} = 0.1$, the
reference-implementation default \citep{labram2024}). The per-model
best-hyperparameter numbers on Stress are presented as observations,
not as controlled comparisons between architectures;
Sec.~\ref{sec:results_paired}'s architecture-independence claim
relies on each architecture at its own best HP on Stress falling
inside a common band, not on cross-architecture ranking.

\subsection{Reproducibility}
All code, training configurations, frozen and fine-tuned checkpoints,
variance-decomposition outputs, permutation-null logs, the pre/post
LaBraM diagnostic, the class-balanced classical audit, and the
contrast-strength anchoring protocol specification will be released
at the project repository upon publication.
